% Generated from validated Article Draft Result. Do not hand-edit; revise the structured draft instead.
\section{ChatGPTは突然の起点ではない――対話UIとcontrol layerが収束した年末}
\label{pkg:dialogue-interface-control}
\noindent\textbf{2022年末のChatGPT research previewは象徴的だったが、その前にはLaMDAのdialogue modeling、helpful/harmless系RLHF、Whisperによるspeech interface、Sparrowやalignment研究、reward-modelの限界に関する検討が積み上がっていた。Constitutional AIまで含めて見ると、対話製品の成立と同時に『どう制御するか』が独立したsystem layerとして前景化したことが分かる。}\autocite{src-2a9e3cfb3148,src-2b132b9f05ef,src-7ef567047bd7,src-bc1a61aec959,src-f4ecc26bd1e5}

ChatGPTの2022年11月公開は、後年の成熟したplatformをそのまま2022年へ持ち込んで読むべき出来事ではない。一次資料上の境界はresearch previewであり、対話UIを通じてinstruction-following modelを広い利用者へ見せたproduct-interface上の転換点として扱うのが適切である。plugin、後年のmultimodal機能、enterprise機能などをこの時点へ逆投影しないことが、2022年を歴史として読む最低条件になる。\autocite{src-bc1a61aec959}


対話interfaceはChatGPTだけから始まったわけではない。LaMDAはdialogue modelを主要な研究対象として示し、Whisperはspeech recognitionを通じて音声をmodel-centered systemへ取り込む別のinterfaceを前景化した。dialogue quality、speech input/output、underlying language-model capabilityは同じ概念ではなく、利用者がmodelと相互作用するsurfaceが複数modalitiesへ広がっていたと見るべきである。\autocite{src-2a9e3cfb3148,src-2b132b9f05ef,src-7ef567047bd7}


対話modelが利用者の指示へ直接応答するほど、alignmentは抽象的な研究agendaだけではなくproduct behaviorを決めるcontrol問題へ近付く。helpful and harmless assistantの研究、OpenAIが示したalignment research approach、DeepMindのSparrowは、それぞれ異なる方法・組織文脈から、望ましい応答と望ましくない応答をどう定義し、model behaviorへ反映するかという問題を扱った。これらを単一の安全方式とみなすのではなく、control layerを構成する複数の試みとして位置付ける。\autocite{src-81cf88ab9c36,src-ce2ca3d4fe8f,src-d7f602cae828}


controlをobjectiveへ落とし込めば問題が解けるわけでもない。reward-model overoptimizationを扱う研究は、proxy objectiveを強く最適化すること自体が別の失敗要因になり得るという境界を示す。Constitutional AIはrule/principleを使ってmodel behaviorを形成する別の方向を提示した。helpfulness/harmlessness、reward signal、明示的なruleという異なるcontrol surfaceが年内に並んだことは、強い対話modelの運用がmodel weightsだけの問題ではないことを示している。\autocite{src-1a032a768300,src-83ce66a102c3,src-ce2ca3d4fe8f,src-f4ecc26bd1e5}


この流れの上にChatGPTを置くと、2022年末は『突然チャットボットが現れた』瞬間ではなく、それまで別々に見えていたinstruction-following、dialogue modeling、speech interface、RLHF、安全規則、reward objectiveの問題が、利用者向け対話UIの周囲で一つのsystemとして見え始めた境界になる。能力、interface、controlを分けて読むことで、2023年以降のproduct competitionへ何が引き継がれたのかも整理しやすくなる。\autocite{src-1a032a768300,src-2a9e3cfb3148,src-2b132b9f05ef,src-7ef567047bd7,src-81cf88ab9c36,src-83ce66a102c3,src-bc1a61aec959,src-ce2ca3d4fe8f,src-d7f602cae828,src-f4ecc26bd1e5}


\begin{claimboundary}
各組織・各論文が示したhelpfulness、harmlessness、安全性、性能、reward挙動は、それぞれの評価条件に帰属する。本稿はそれらを独立に再現した安全性ランキングへ換算せず、research agenda、deployed interface、rule/reward-based controlの状態を分けて記述する。\autocite{src-1a032a768300,src-83ce66a102c3,src-f4ecc26bd1e5}
\end{claimboundary}
